\section{Introduction}
\label{sec:introduction}
RDF knowledge graphs (KGs) have become a standard framework for structuring open-domain information, as exemplified by knowledge bases such as DBpedia and Wikidata \cite{auer2007dbpedia, vrandecic2014wikidata}. Simultaneously, their application has expanded to capture complex, domain-specific knowledge in areas ranging from soil carbon and climate modeling to medicine and sociology \cite{frink_proto-okn}. SPARQL has become the standard query language for RDF knowledge graphs, with hundreds of queryable endpoints emerging since its inception \cite{ali2022survey}. Owing to the technical barrier of composing SPARQL queries, Knowledge Graph Question Answering (KGQA) systems have become essential for democratizing access to knowledge graphs, enabling researchers, domain-experts, and ordinary users to seamlessly reason about such knowledge by mapping natural language questions to SPARQL queries that can be executed against the graph to provide accurate answers. 
% Crucially, these systems guard against model hallucinations in an era where large language models are the primary front-end for information access.
\todo{discussing KGQA in the context of Ai.}\todo[color=green]{the users are not just researchers. it could be domain experts, non experts, and agentic systems too.}\todo[color=green]{make it more convincing and explicit that SPARQL is still crucial.}

Text-SPARQL example sets play a central role for KGQA systems.\todo[color=cyan]{More systematic exploration of the distinction between 3 natural language query formats: keywords, questions, and commands.} They serve as in-context demonstrations that steer LLM-based query generation \cite{wikisp2023,zahera2024generating, dabramo-etal-2025-investigating,emonet2024llm}, as training corpora for fine-tuning translation models \cite{sgpt2022generative,neural2022machine}, and as standardized benchmarks for measuring system correctness \cite{dubey2019lcquad, usbeck2019qald,usbeck2021qald}. Beyond training and evaluation, these example sets are often used as executable documentation, helping human practitioners better navigate unfamiliar knowledge graph structures and query idioms \cite{uniprot_sparql_examples, wikidata_sparql_examples}.

Constructing high-quality Text-SPARQL example sets for a given knowledge graph is a costly and time-consuming process because it requires three distinct areas of expertise: (i) fluency in the specific target domain, (ii) proficiency in formal SPARQL query construction, and (iii) a deep familiarity with the knowledge graph's ontology, including its semantics, structure, and vocabulary. Existing curated Text-SPARQL benchmarks concentrate on a handful of prominent open-domain knowledge graphs, principally DBpedia, Wikidata, and Freebase, and rely on manual question authoring, template-based construction, and contributions from the open-source community \cite{usbeck2019qald,usbeck2021qald,trivedi2017lc,dubey2019lcquad,wikisp2023,liu2024spinach}.\todo[color=green]{Make it clear that QALD and so on are such existing benchmarks on open-domain KGs.}\todo[color=green]{Unclear why we discuss each benchmark's creation} Domain-specific knowledge graphs, which encode specialized knowledge in fields such as biomedicine \cite{kgs2025biomedical}, geoscience \cite{sockg2025}, and climate models \cite{adamu2025climatepub4kg}, lack equivalent community infrastructure and consequently have few or no Text-SPARQL resources \cite{frink_proto-okn}. This scarcity has confined the majority of KGQA research \cite{liu2024spinach,zahera2024generating,sgpt2022generative,kacupaj2023sparqlgen,banerjee2022modern,diallo2025frase,sun2022think,wikisp2023} to the open-domain settings where Text-SPARQL example sets already exist, leaving domain-specific settings largely unaddressed, and depriving domain researchers of the automated QA interfaces necessary to seamlessly reason over specialized knowledge embedded in domain-specific knowledge graphs.

To address this gap, there is a need for an automated KG-agnostic methodology that can generate reliable, diverse and comprehensive Text-SPARQL pairs for arbitrary knowledge graphs, which to the best of our knowledge no prior work has yet addressed. To bridge this gap, we present QueryBloom, an automated, KG-agnostic methodology that can synthesize high-quality Text-SPARQL example sets from any RDFS-based knowledge graph given only SPARQL endpoint access.\todo[color=green]{Need to articulate the motivations behind the methodology design before you detail the methods.} The approach is motivated by the two following observations: (i), it is possible to discover arbitrary subgraph patterns for any given RDFS-based knowledge graph using KG-agnostic RDFS-based exploration queries (Subgraph Discoverability) and (ii), every SPARQL query implicitly specifies a subgraph pattern over the knowledge graph (SPARQL-Subgraph Duality). QueryBloom inverts the duality by using an LLM agentic pipeline to generate Text-SPARQL queries grounded on subgraph instances and patterns automatically discovered from the target knowledge graph.

Concretely, the QueryBloom pipeline proceeds in four stages: (1) extracting one-hop subgraph type patterns encoding which properties connect which types of entities; (2) growing a set of structurally diverse multi-hop subgraphs by iteratively extending seed patterns through adjacency-guided attachment; (3) enriching each subgraph with concrete entity instances and attribute values; and (4) employing an agentic pipeline that filters non-meaningful subgraphs, jointly generates natural language questions and SPARQL queries, augments the pairs to introduce meaningful complexity, iteratively verifies each pair against the endpoint, and finally refines phrasing for the generated questions.

\todo{Explain the evaluation, summarize the results, and other aspects before the ``summary''.}

In summary, our key contributions are as follows:\todo[color=green]{can have more bullet points.}
\begin{itemize}
	\item QueryBloom, an LLM-driven methodology that synthesizes diverse, accurate and, comprehensive Text-SPARQL examples for a target knowledge graph given only its SPARQL endpoint. It tackles the example scarcity challenge for domain-specific knowledge graphs and facilitates
 the development of their KGQA systems.
	\item A comprehensive evaluation framework demonstrating that QueryBloom-generated datasets match or exceed the diversity of established human-authored benchmarks across query structural complexity, semantic linguistic variety, and holistic schema coverage.
	\item Three novel Text-to-SPARQL datasets across diverse domain scopes: open-domain (\textsc{DBpedia}), broad-domain (\textsc{DBLP}), and domain-specific (\textsc{SOCKG}).
\end{itemize}